\label{sec:workflow}

Table~\ref{tab:workflow} presents a practical workflow for all three adoption
pathways, identifying the responsible actor, the \texttt{redist} command,
and the oversight mechanism at each stage.

\begin{table}[h]
\centering
\caption{Practical Workflow: Three Adoption Pathways}
\label{tab:workflow}
\begin{threeparttable}
\begin{tabular}{p{2cm}p{3.5cm}p{3.5cm}p{3.5cm}}
\toprule
Stage & Commission (AZ/CA) & Statute & Court Order (NC) \\
\midrule
\textbf{Data prep} &
Commission staff; \newline
\texttt{bisect fetch} &
State GIS office; \newline
\texttt{bisect fetch} &
Special master staff; \newline
\texttt{bisect fetch} \\
\addlinespace
\textbf{Initial run} &
Staff run \texttt{bisect state}; \newline
Commission reviews output &
Designated agency runs \texttt{bisect state}; \newline
Legislative committee reviews &
Special master runs \texttt{bisect state}; \newline
Court receives report \\
\addlinespace
\textbf{Analysis} &
\texttt{bisect analyze} \newline compactness + demographic &
\texttt{bisect analyze} \newline all types &
\texttt{bisect analyze} \newline political + demographic \\
\addlinespace
\textbf{Public comment} &
60-day public hearing period; \newline Commission adjustments logged &
Legislative hearing; \newline Public comment period &
14-day objection period; \newline Parties file briefs \\
\addlinespace
\textbf{Adjustment} &
Commission vote on adjustments; \newline documented against algorithm baseline &
Agency or commission adjusts; \newline legislative approval &
Special master proposes adjustments; \newline court approves \\
\addlinespace
\textbf{Final adoption} &
Commission vote &
Legislative vote or commission order &
Court order \\
\addlinespace
\textbf{Review} &
State court (if challenged) &
State court (if challenged) &
Higher court (appeal) \\
\bottomrule
\end{tabular}
\begin{tablenotes}
\small
\item Note: \texttt{bisect fetch} downloads and validates census data.
  \texttt{bisect state} runs the partitioning algorithm.
  \texttt{bisect analyze} generates criteria-compliance reports.
  Commands follow the \texttt{bisect} CLI reference~\citep{redistCliQuickstart}.
\end{tablenotes}
\end{threeparttable}
\end{table}

\subsection{Key Differences Across Pathways}

The three pathways differ primarily in three dimensions.

\textbf{Who runs the algorithm.}  In the commission pathway, commission staff
or a designated technical contractor runs the algorithm under the commission's
supervision.  In the statute pathway, a designated state agency (e.g., the
state GIS office or a redistricting board) runs the algorithm.  In the
court-order pathway, the special master or the master's technical support team
runs the algorithm.  In all three cases, the algorithm is run openly, with
full documentation of inputs, parameters, and outputs.

\textbf{How adjustments are made.}  Commission pathways allow the commission
to adjust the algorithm output based on community-of-interest testimony.
Statute pathways may delegate adjustment authority to an agency or reserve it
to the legislature.  Court-order pathways confine adjustments to those
supported by specific constitutional findings, and the special master's
adjustment proposals are subject to judicial review.

\textbf{What counts as a successful challenge.}  In the commission and statute
pathways, successful challenges require showing that the commission or agency
exceeded its authority or violated a criterion.  In the court-order pathway,
the challenge must show a constitutional defect in the output---a higher bar
that makes the court-order pathway the most legally stable of the three.
